\begin{center}\begin{tikzpicture}[font=\small,>=stealth,line join=round]
\draw[gray!40,rounded corners=3pt] (0,0) rectangle (14.2,8.3);\node at (7.1,7.700000000000001) {横纵顺序同时反转，中心对称两侧同向单调};
\begin{scope}[shift={(1.2428571428571429,2.8)},x=1.4285714285714288cm,y=0.575cm]\draw[->,gray] (-0.1,0)--(4.1,0) node[right]{$x$};\draw[->,gray] (0,-2)--(0,6) node[above]{$y$};\draw[blue,thick,solid,domain=0.5:3.5,samples=180] plot(\x,{(\x-2)^3+2});\fill[orange!85!black] (0.5,-1.375) circle(2.2pt);\fill[orange!85!black] (1,1) circle(2.2pt);\fill[orange!85!black] (3,3) circle(2.2pt);\fill[orange!85!black] (3.5,5.375) circle(2.2pt);\end{scope}\node[anchor=west,align=left,text width=5.0cm] at (8.4,4.0) {示例 $f=(x-2)^3+2$\\[9pt]中心 $(2,2)$\\[9pt]左侧 $0.5<1$\\[9pt]右侧对应 $3.5>3$\\[9pt]$f(0.5)<f(1)$ 仍成立};
\node[text width=12.6cm,align=center] at (7.1,.65) {图取例6的一个具体函数；证明用比较式，适用于所有满足题设的函数。};
\end{tikzpicture}\end{center}
